\documentclass[../../../include/open-logic-section]{subfiles}

\begin{document}

\olfileid{sth}{card-arithmetic}{opps}
\olsection{મૂળભૂત ક્રિયાઓની વ્યાખ્યા}

આપણને ક્રમની નોંધ રાખવાની જરૂર નથી, તેથી
ગણાંકોનું અંકગણિત ક્રમવાચક અંકગણિત કરતાં વ્યાખ્યાયિત
કરવામાં સહેલું છે. સરવાળો, ગુણાકાર અને ઘાતક્રિયા
ત્રણેયને સાથે વ્યાખ્યાયિત કરીશું.

\begin{defn}
$\cardfont{a}$ અને $\cardfont{b}$ ગણાંકો હોય ત્યારે:
\begin{align*}
	\cardfont{a} \cardplus \cardfont{b} &\defis 
	\card{\cardfont{a} \disjointsum \cardfont{b}}\\
	\cardfont{a} \cardtimes \cardfont{b} &\defis 
	\card{\cardfont{a} \times \cardfont{b}}\\
	\cardexpo{\cardfont{a}}{\cardfont{b}} &\defis 
	\card{\funfromto{\cardfont{b}}{\cardfont{a}}}
\end{align*}
અહીં $\funfromto{X}{Y} = \Setabs{f}{f \text{ એ વિધેય છે } X \to Y}$.
(કોઈ પણ ગણો $X$ અને $Y$ માટે
$\funfromto{X}{Y}$ અસ્તિત્વમાં છે તે બતાવવું સહેલું
છે; તેને સ્વાધ્યાય તરીકે છોડીએ છીએ.)
\end{defn}

\begin{prob}
કોઈ પણ ગણો $X$ અને~$Y$ માટે
$\funfromto{X}{Y}$ અસ્તિત્વમાં છે તે $\Zminus$માં
સાબિત કરો. $\ZF$માં કામ કરતાં, $\setrank{X}$ અને
$\setrank{Y}$ પરથી $\setrank{\funfromto{X}{Y}}$
ગણો; \olref[sth][ord-arithmetic][using-addition]{rankcomputation}
જેવી રીતે.
\end{prob}

આ વ્યાખ્યા થોડી સમજાવીએ. સરવાળા માટે તે
\olref[ord-arithmetic][add]{defdissum}માં
વ્યાખ્યાયિત વિચ્છિન્ન સરવાળો $\disjointsum$ વાપરે
છે; સાન્ત કિસ્સામાં તે યોગ્ય પરિણામ આપે છે તે
સરળતાથી દેખાય છે. ગુણાકાર માટે
\olref[sfr][set][pai]{cardnmprod} કહે છે કે $A$માં
$n$ ઘટકો અને $B$માં $m$ ઘટકો હોય, તો
$A \times B$માં $n \cdot m$ ઘટકો હોય. તેથી
આપણી વ્યાખ્યા એ જ વિચારને અનંતાતીત ગુણાકાર સુધી
સામાન્ય બનાવે છે. ઘાતક્રિયા માટે પણ એવું જ છે:
સાન્ત કિસ્સાનો વિચાર અનંતાતીત કિસ્સા સુધી લઈ જઈએ
છીએ. ખરેખર કેટલીક રીતે અનંતાતીત ગણાંક-અંકગણિત
અનંતાતીત ક્રમવાચક અંકગણિત કરતાં ``સામાન્ય''
અંકગણિત જેવું વધુ લાગે છે:

\begin{prop}\ollabel{cardplustimescommute}
$\cardplus$ અને $\cardtimes$ અદલબદલીય અને સહયોગી છે.
\end{prop}

\begin{proof}
અદલબદલી માટે
\olref[cardinals][cardsasords]{lem:CardinalsBehaveRight}
મુજબ $\cardeq{(\cardfont{a} \disjointsum
\cardfont{b})}{(\cardfont{b} \disjointsum \cardfont{a})}$
અને $\cardeq{(\cardfont{a} \times
\cardfont{b})}{(\cardfont{b} \times \cardfont{a})}$
નોંધવું પૂરતું છે. સહયોગિતાનો પુરાવો સ્વાધ્યાય
માટે છોડીએ છીએ.
\end{proof}

\begin{prob}
$\cardplus$ અને $\cardtimes$ સહયોગી છે તે સાબિત કરો.
\end{prob}

\begin{prop}
$A$ અનંત હોય જો અને માત્ર જો
$\card{A} \cardplus 1 = 1 \cardplus \card{A} = \card{A}$.
\end{prop}

\begin{proof}
\olref[cardinals][classing]{generalinfinitycharacter}ની
જેમ, \olref[ord-arithmetic][using-addition]{ordinfinitycharacter}
અને \olref[cardinals][cardsasords]{lem:CardinalsBehaveRight}
પરથી મળે છે.
\end{proof}

આથી સમજાય છે કે ક્રમવાચક સરવાળા-ગુણાકાર અને
ગણાંક સરવાળા-ગુણાકાર માટે જુદાં સંકેતો શા માટે
જોઈએ: એ ખરેખર \emph{જુદી} ક્રિયાઓ છે. આગળનાં
બે પરિણામો બતાવે છે કે ક્રમવાચક અને ગણાંક ઘાતક્રિયા
પણ જુદી છે. (યાદ કરો કે
\olref[z][infinity-again]{defnomega}થી $2 = \{0,
1\}$ મળે છે):

\begin{lem}\ollabel{lem:SizePowerset2Exp}
કોઈ પણ $A$ માટે $\card{\Pow{A}} =
\cardexpo{2}{\card{A}}$.
\end{lem}

\begin{proof}
દરેક ઉપગણ $B \subseteq A$ માટે
$\chi_B \in \funfromto{A}{2}$ આ રીતે લો:
\begin{align*}
	\chi_{B}(x) &\defis
	\begin{cases}
		1 & \text{જો }x\in B\\
		0 & \text{અન્યથા.}
	\end{cases}
\end{align*}
હવે $f(B) = \chi_B$ લો; આથી
!!a{bijection} $f \colon \Pow{A} \to
\funfromto{A}{2}$ વ્યાખ્યાયિત થાય છે.
તેથી $\cardeq{\Pow{A}}{\funfromto{A}{2}}$.
આથી $\cardeq{\Pow{A}}{\funfromto{\card{A}}{2}}$,
અને તેથી $\card{\Pow{A}} =
\card{\funfromto{\card{A}}{2}} = 2^{\card{A}}$.
\end{proof}

આ ટૂંકો પુરાવો એક રીતે
\olref[sfr][siz][red-alt]{sec}ની ચર્ચાને પણ સમાવે છે.
ત્યાં આપણે $\Pow{\omega}$ની અગણનીયતાને અનંત દ્વિઆધારી
અનુક્રમોના ગણ $\Bin^\omega$ની અગણનીયતામાં કેવી રીતે
``ઘટાડી'' શકાય તે બતાવ્યું હતું. હકીકતમાં
$\Bin^{\omega}$ એટલે $\funfromto{\omega}{2}$ જ;
અને ઉપરનો પુરાવો બતાવે છે કે
\olref[sfr][siz][red-alt]{sec}ની દલીલમાં~$\omega$
ને બદલે કોઈ પણ ગણ~$A$ લઈ શકીએ. પરિણામમાંથી
ગણાંક ઘાતક્રિયા અંગે એક ટૂંકું તથ્ય પણ મળે છે:

\begin{cor}\ollabel{cantorcor}
કોઈ પણ ગણાંક~$\cardfont{a}$ માટે
$\cardfont{a} < \cardexpo{2}{\cardfont{a}}$.
\end{cor}

\begin{proof}
કૅન્ટરના પ્રમેય
(\olref[sfr][siz][car]{thm:cantor}) અને
\olref{lem:SizePowerset2Exp} પરથી.
\end{proof}
\noindent
તેથી $\omega < \cardexpo{2}{\omega}$.
પરંતુ ધ્યાન રાખો: આ પરિણામ \emph{ગણાંક}
ઘાતક્રિયા વિશે છે. તેની સરખામણી \emph{ક્રમવાચક}
ઘાતક્રિયા સાથે કરવી જોઈએ, કારણ કે ત્યાં
$\omega = \ordexpo{2}{\omega}$ છે
(\olref[ord-arithmetic][expo]{sec} જુઓ).

ગણાંક ઘાતક્રિયાની વાત ચાલે છે ત્યારે
$\Real$ કઈ ``રીતે'' !!{nonenumerable} છે તે
વધુ ચોક્કસ કહી શકીએ.

\begin{thm}\ollabel{continuumis2aleph0}
$\card{\Real} = \cardexpo{2}{\omega}$
\end{thm}

\begin{proof}[પુરાવાની રૂપરેખા]
આ સાબિત કરવાની ઘણી રીતો છે. સૌથી સીધી રીત એ
છે કે $\cardle{\Pow{\omega}}{\Real}$ અને
$\cardle{\Real}{\Pow{\omega}}$ બતાવીએ; પછી
શ્રેડર--બર્નસ્ટાઇનથી $\cardeq{\Real}{\Pow{\omega}}$
મળે અને \olref[card-arithmetic][opps]{lem:SizePowerset2Exp}
થી $\card{\Real} = \cardexpo{2}{\omega}$ મળે.
અંતઃક્ષેપો $f \colon \Pow{\omega} \to \Real$
અને $g \colon \Real \to \Pow{\omega}$ વ્યાખ્યાયિત
કરવાનું (ઉપયોગી) સ્વાધ્યાય માટે છોડીએ છીએ.
\end{proof}

\begin{prob}
$\cardle{\Pow{\omega}}{\Real}$ અને
$\cardle{\Real}{\Pow{\omega}}$ બતાવીને
\olref[sth][card-arithmetic][opps]{continuumis2aleph0}નો
પુરાવો પૂર્ણ કરો.
\end{prob}

\end{document}
